% DERIVED from formal_layer.json — read-only, do not edit.
\begin{definitionv}{synth_4}[Fixed public regularity radius]
The public regularity radius \(L\) is the real constant
\[
L = 3.
\]
\end{definitionv}

\begin{definitionv}{synth_9}[Observed records and ordered datasets]
The observed-record space is
\[
\mathcal O=[0,1]\times\{0,1\}^2,
\]
equipped with its Borel sigma-algebra, with a record written as \(z=(x,a,y)\). For a nonnegative integer \(n\), define \(\mathcal O^n\) to be its ordered \(n\)-fold product, equipped with the product sigma-algebra. A dataset is an element
\[
D=(z_1,\ldots,z_n)\in\mathcal O^n.
\]
For a random dataset \(D\), write \(\mathcal L(D)\) for its probability law on \(\mathcal O^n\); explicitly,
\[
\mathcal L(D)(E)=\Pr\{D\in E\}
\]
for every measurable \(E\subseteq\mathcal O^n\). For \(n=0\), the product consists of the empty dataset.
\end{definitionv}

\begin{definitionv}{synth_3}[Observed binary outcome coordinate]
The observed binary outcome \(Y\) is the final coordinate of a full causal record, whose coordinates are, in order, a covariate, a binary treatment indicator, an ordered pair of binary potential outcomes, and a binary observed outcome.
\end{definitionv}

\begin{definitionv}{synth_2}[Binary treatment coordinate]
The binary treatment indicator \(A\) maps each full causal record to its treatment coordinate, taking values in \(\{0,1\}\). A full causal record is ordered as covariate, binary treatment, ordered pair of binary potential outcomes, and binary observed outcome.
\end{definitionv}

\begin{definitionv}{synth_1}[Scalar covariate coordinate]
The scalar covariate \(X\) is the first coordinate of the causal record, with values in the closed unit interval:
\[
X \in [0,1].
\]
\end{definitionv}

\begin{definitionv}{synth_8}[Causal coordinates and marginal laws]
A causal law \(P\) is a Borel probability law of
\[
(X,A,Y^{\mathrm{pot}},Y)
\in[0,1]\times\{0,1\}\times\{0,1\}^2\times\{0,1\},
\]
where \(X\) is the covariate, \(A\) is the treatment indicator, \(Y\) is the observed outcome, and
\[
Y^{\mathrm{pot}}=(Y(0),Y(1))
\]
is the ordered pair of control and treatment potential outcomes. Define the observational marginal and the covariate marginal by
\[
P^{\mathrm{obs}}=\mathcal L_P(X,A,Y),
\qquad
P_X=\mathcal L_P(X).
\]
Here \(\mathcal L_P\) denotes the probability law of the indicated coordinates under \(P\).
\end{definitionv}

\begin{assumptionv}{ass:iid}[Independent and identically distributed observations]
The dataset \(D\) has distribution
\[
\mathcal L(D)=(P^{\mathrm{obs}})^{\otimes n}.
\]
\end{assumptionv}

\begin{assumptionv}{ass:consistency}[Consistency of observed outcomes]
The observed outcome \(Y\) equals the potential outcome \(Y(A)\) selected by the realized treatment with probability one:
\[
P\{Y=Y(A)\}=1.
\]
\end{assumptionv}

\begin{assumptionv}{ass:exchangeability}[Conditional exchangeability of treatment]
The potential-outcome pair \(Y^{\mathrm{pot}}\) is conditionally independent of the treatment indicator \(A\) given the covariate \(X\) under the causal law \(P\):
\[
Y^{\mathrm{pot}}\perp\!\!\!\perp A\mid X\quad[P].
\]
\end{assumptionv}

\begin{definitionv}{synth_10}[Selected design density and observational regressions]
For a causal law \(P\) whose covariate marginal \(P_X\) is absolutely continuous with respect to Lebesgue measure \(dx\) on \([0,1]\), let \(f_P\) be the selected measurable version of its density:
\[
f_P=\frac{dP_X}{dx},
\qquad
P_X(B)=\int_B f_P(x)\,dx
\]
for every Borel set \(B\subseteq[0,1]\). Let \(e_P\) and \(\mu_{0,P}\) be the selected versions of the propensity regression and the control outcome regression, respectively:
\[
e_P(x)=P(A=1\mid X=x),
\qquad
\mu_{0,P}(x)=E_P[Y\mid X=x,A=0],
\qquad x\in[0,1].
\]
These conditional expressions denote the selected regression functions on \([0,1]\). The model's overlap and smoothness conditions apply to these selected versions, and its density bounds apply to \(f_P\) Lebesgue-almost everywhere.
\end{definitionv}

\begin{assumptionv}{ass:density}[Bounded covariate density]
The covariate density \(f_P\) satisfies
\[
\frac12\le f_P(x)\le\frac32
\quad\text{for Lebesgue-almost every }x\in[0,1].
\]
\end{assumptionv}

\begin{assumptionv}{ass:overlap}[Uniform treatment overlap]
The conditional treatment probability \(e_P\) satisfies
\[
\frac14\le e_P(x)\le\frac34
\qquad\text{for every }x\in[0,1].
\]
\end{assumptionv}

\begin{assumptionv}{ass:propensity-holder}[Hölder continuity of propensity]
The conditional treatment probability \(e_P\) satisfies the Hölder bound with common smoothness radius \(L\):
\[
|e_P(x)-e_P(x^{\prime})|
\le L|x-x^{\prime}|^{1/10}
\qquad\text{for every }x,x^{\prime}\in[0,1].
\]
\end{assumptionv}

\begin{assumptionv}{ass:control-holder}[Hölder continuity of control means]
For every \(x,x'\in[0,1]\),
\[
|\mu_{0,P}(x)-\mu_{0,P}(x')|
\le L|x-x'|^{1/10}.
\]
\end{assumptionv}

\begin{assumptionv}{ass:contrast-lipschitz}[Lipschitz continuity of contrasts]
For every \(x,x'\in[0,1]\),
\[
|\tau_P(x)-\tau_P(x')|
\le L|x-x'|.
\]
\end{assumptionv}

\begin{definitionv}{synth_16}[Borel probability laws]
For a space \(\mathsf X\) equipped with its Borel \(\sigma\)-algebra \(\mathcal B(\mathsf X)\), define
\[
\operatorname{Prob}(\mathsf X)
=
\{P:P\text{ is a probability measure on }(\mathsf X,\mathcal B(\mathsf X))\}.
\]
\end{definitionv}

\begin{definitionv}{def:complete-model}[Causal law class \(\mathcal P\)]
\(\mathcal P\) is the class of causal probability laws
\[
P\in\operatorname{Prob}
\bigl([0,1]\times\{0,1\}\times\{0,1\}^2\times\{0,1\}\bigr)
\]
satisfying all of the following conditions:
\begin{itemize}
\item \textbf{(Consistency.)} \cref{obj:ass:consistency}.
\item \textbf{(Exchangeability.)} \cref{obj:ass:exchangeability}.
\item \textbf{(Covariate density.)} \cref{obj:ass:density}.
\item \textbf{(Overlap.)} \cref{obj:ass:overlap}.
\item \textbf{(Propensity smoothness.)} \cref{obj:ass:propensity-holder}.
\item \textbf{(Control-mean smoothness.)} \cref{obj:ass:control-holder}.
\item \textbf{(Contrast smoothness.)} \cref{obj:ass:contrast-lipschitz}.
\end{itemize}
The function conditions apply to the selected versions of the observational regressions. The covariate density is an unknown measurable function. Both arm means range over their full binary-outcome domain, subject to the stated conditions. Membership in \(\mathcal P\) is determined by these properties of the causal law.
\end{definitionv}

\begin{assumptionv}{ass:pure-privacy}[Pure replacement differential privacy]
For every pair of datasets \(D,D'\in\mathcal O^n\) with replacement distance \(d_{\mathrm H}(D,D')=1\), and every measurable output event \(E\in\mathcal B(\mathcal B)\),
\[
M(D,E)\le \exp(\epsilon)M(D',E).
\]
\end{assumptionv}

\begin{definitionv}{def:private-kernels}[Private release class \(\mathfrak M_{n,\epsilon}(\mathcal B)\)]
\(\mathfrak M_{n,\epsilon}(\mathcal B)\) is the class of all everywhere-defined Markov kernels
\[
M:\mathcal O^n\rightsquigarrow\mathcal B
\]
satisfying \cref{obj:ass:pure-privacy}. Here \(\mathcal B\) is a standard Borel output space, and \(M(D,E)\) is the conditional probability of an output in the Borel event \(E\) given dataset \(D\).
\end{definitionv}

\begin{definitionv}{def:risk-criterion}[Minimax absolute-error risk \(R_{n,\epsilon}\)]
\[
R_{n,\epsilon}
=
\inf_{T\in\mathfrak M_{n,\epsilon}(\mathbb R)}
\sup_{P\in\mathcal P}
E_{(P^{\mathrm{obs}})^{\otimes n},T}
\bigl|T(D)-\theta(P)\bigr|.
\]
Here \(T(D)\) is the scalar output of the randomized estimator, \(\theta(P)\) is the pointwise treatment-effect target, and the expectation integrates both the observed sample with law \((P^{\mathrm{obs}})^{\otimes n}\) and the estimator randomization.
\end{definitionv}

\begin{definitionv}{synth_13}[Measurable interval output space]
Define
\[
\mathcal I=\{I\subseteq\mathbb R:I\text{ is a connected Borel set}\}.
\]
Code the empty set by a distinguished symbol \(\varnothing\). Code each nonempty interval by its lower and upper endpoints \(\ell,u\) and endpoint-inclusion flags \(\eta_\ell,\eta_u\in\{0,1\}\), where
\[
\ell\in\mathbb R\cup\{-\infty\},\qquad
u\in\mathbb R\cup\{+\infty\},\qquad \ell\le u.
\]
For \(\ell<u\), the code represents
\[
(\ell,u)\cap\mathbb R
\;\cup\;
\begin{cases}\{\ell\},&\ell\in\mathbb R,\ \eta_\ell=1,\\ \varnothing,&\text{otherwise},\end{cases}
\;\cup\;
\begin{cases}\{u\},&u\in\mathbb R,\ \eta_u=1,\\ \varnothing,&\text{otherwise}.\end{cases}
\]
An infinite endpoint has inclusion flag zero. For \(\ell=u\), require \(\ell=u\in\mathbb R\) and \(\eta_\ell=\eta_u=1\), so the code represents the singleton \(\{\ell\}\). Equip the valid nonempty codes with the Borel sigma-algebra inherited from
\[
(\mathbb R\cup\{-\infty\})\times
(\mathbb R\cup\{+\infty\})\times\{0,1\}^2,
\]
using the usual extended-real Borel structures and discrete flags, and adjoin the empty-set code as an isolated point. Transport this sigma-algebra to \(\mathcal I\) through the coding bijection. This makes \(\mathcal I\) a standard Borel measurable space, including empty, singleton, bounded, and unbounded interval outputs.
\end{definitionv}

\begin{definitionv}{def:interval-criterion}[Honest expected-length criterion \(H_{n,\epsilon}\)]
The minimax expected length under uniform \(90\%\) coverage is
\[
H_{n,\epsilon}
=
\inf_{\substack{
I\in\mathfrak M_{n,\epsilon}(\mathcal I)\\
\inf_{P\in\mathcal P}
\Pr_{(P^{\mathrm{obs}})^{\otimes n},I}
\{\theta(P)\in I(D)\}\ge 9/10
}}
\ \sup_{P\in\mathcal P}
E_{(P^{\mathrm{obs}})^{\otimes n},I}
\operatorname{Leb}(I(D)).
\]
Here \(\mathcal I\) is the interval output space, \(I(D)\) is the released connected interval, and \(\operatorname{Leb}(I(D))\) is its Lebesgue length. The probability and expectation integrate both the observed sample and the procedure's randomization.
\end{definitionv}

\begin{definitionv}{synth_12}[Clipping to a closed interval]
For real numbers \(a\le b\) and \(u\in\mathbb R\), define
\[
\operatorname{clip}_{[a,b]}(u)=\min\{b,\max\{a,u\}\}.
\]
\end{definitionv}

\begin{definitionv}{synth_11}[Public local cells]
Let \(x_0=1/2\), let \(h\in(0,1/4]\) be the public localization radius, and let \(k\ge2\) be the integer public cell count. Set \(\delta=2h/k\). The public local window is \([x_0-h,x_0+h]\subseteq[0,1]\), with indexed partition \((C_j)_{j=1}^k\) defined by
\[
C_j=
\begin{cases}
[x_0-h+(j-1)\delta,\ x_0-h+j\delta),&1\le j<k,\\
[x_0-h+(k-1)\delta,\ x_0+h],&j=k.
\end{cases}
\]
Public cell \(j\) means \(C_j\). A record \((X_i,A_i,Y_i)\) belongs to public cell \(j\) precisely when \(X_i\in C_j\); its count is
\[
N_j(D)=\sum_{i=1}^n\mathbf1\{X_i\in C_j\}.
\]
Records are counted with their dataset multiplicities.
\end{definitionv}

\begin{algorithmv}{def:private-ratio}[Private ratio \(\widehat T_{h,k}\) and interval \(\widehat I_{h,k}\)]
The inputs are an ordered dataset \(D\in\mathcal O^n\), the public localization radius \(h\), the public cell count \(k\), the privacy parameter \(\epsilon\), the public occupancy scale \(\mathcal A\), and the public error envelope \(V\).
\begin{enumerate}
\item Set the cell width and positive denominator floor to
\[
\delta=\frac{2h}{k},
\qquad
d_0=\frac{3\mathcal A}{64}>0.
\]
For public cell \(j\), let its records and record count be
\[
D_j=\{(X_i,A_i,Y_i):i\text{ indexes a record of }D
\text{ in public cell }j\},
\qquad
s_j=|D_j|,
\]
where records are counted with their dataset multiplicities.

\item Form the numerator and denominator statistics by summing the cell contributions:
\[
S_N(D)=
\sum_j
\begin{cases}
0,&s_j<2,\\[2pt]
\displaystyle
\frac{2\left\{
s_j\sum_{(X,A,Y)\in D_j}AY
-\left(\sum_{(X,A,Y)\in D_j}A\right)
 \left(\sum_{(X,A,Y)\in D_j}Y\right)
\right\}}{s_j-1},&s_j\ge2,
\end{cases}
\]
\[
S_D(D)=
\sum_j
\begin{cases}
0,&s_j<2,\\[2pt]
\displaystyle
\frac{2\left\{
s_j\sum_{(X,A,Y)\in D_j}A
-\left(\sum_{(X,A,Y)\in D_j}A\right)^2
\right\}}{s_j-1},&s_j\ge2.
\end{cases}
\]
Only occupied cells need to be included in these sums.

\item Draw independent Laplace noises once:
\[
\xi_N,\xi_D\ \stackrel{\mathrm{ind}}{\sim}\
\operatorname{Laplace}\!\left(0,\frac{12}{\epsilon}\right).
\]
Release the joint statistic
\[
\widetilde S
=
(\widetilde S_N,\widetilde S_D)
=
\bigl(S_N(D)+\xi_N,\ S_D(D)+\xi_D\bigr).
\]

\item Output the following postprocessing of this single release:
\[
\widehat T_{h,k}
=
\operatorname{clip}_{[-1,1]}
\left\{
\frac{\widetilde S_N}
{\max\{d_0,\widetilde S_D\}}
\right\},
\]
\[
\widehat I_{h,k}
=
\left[
\max\left\{-1,\widehat T_{h,k}-\sqrt{10V(h,\delta)}\right\},
\min\left\{1,\widehat T_{h,k}+\sqrt{10V(h,\delta)}\right\}
\right].
\]
The map is defined on every \(D\in\mathcal O^n\).
\end{enumerate}
\end{algorithmv}

\begin{algorithmv}{def:public-tuning}[Public tuning for \(\widehat T^*\) and \(\widehat I^*\)]
The inputs are the dataset, the sample size \(n\), the privacy parameter \(\epsilon\), and the numerical benchmark \(r(n,\epsilon)\).
\begin{enumerate}
\item If \(r(n,\epsilon)\ge1/8\), output
\[
\widehat T^*=0,
\qquad
\widehat I^*=[-1,1].
\]

\item If \(r(n,\epsilon)<1/8\), choose the public tuning parameters
\[
h=2^{-\lceil\log_2(1/r(n,\epsilon))\rceil},
\qquad
k=2h^{-4},
\qquad
\delta=\frac{2h}{k}=h^5,
\]
and output the local private ratio and interval:
\[
(\widehat T^*,\widehat I^*)
=
(\widehat T_{h,k},\widehat I_{h,k}).
\]
\end{enumerate}
This defines a single sample map whose public tuning depends only on \(n\) and \(\epsilon\).
\end{algorithmv}

\begin{definitionv}{synth_17}[Interval hull]
For a nonempty subset \(S\subseteq\mathbb R\), define its interval hull as the smallest connected interval containing \(S\):
\[
\operatorname{hull}(S)
=
\{x\in\mathbb R:\text{there exist }a,b\in S\text{ with }a\le x\le b\}.
\]
\end{definitionv}

\begin{algorithmv}{def:interval-handle}[Interval inversion $I^*_{n,\epsilon}$]
The inputs are the sample size \(n\), privacy parameter \(\epsilon\), benchmark \(r(n,\epsilon)\), public tuning parameters \(h,k,\delta\), denominator floor \(d_0\), and mean-squared-error envelope \(V(h,\delta)\). The envelope incorporates the bias envelope \(b\), occupancy scale \(\mathcal A\), and Laplace second moment.
\begin{enumerate}
\item If \(r(n,\epsilon)\ge 1/8\), accept every candidate \(\vartheta\in[-1,1]\) and output
\[
I^*_{n,\epsilon}=\widehat I^*=[-1,1].
\]
Otherwise, proceed to the following steps.
\item Compute the single joint noisy release \(\widetilde S\), with numerator coordinate \(\widetilde S_N\) and denominator coordinate \(\widetilde S_D\), using the public tuning parameters. Define
\[
B_D=\max\{d_0,\widetilde S_D\},\qquad
B_N=\operatorname{clip}_{[-B_D,B_D]}(\widetilde S_N),\qquad
\rho=\sqrt{10V(h,\delta)}.
\]
\item Accept a candidate \(\vartheta\in[-1,1]\) precisely when
\[
|B_N-\vartheta B_D|\le \rho B_D.
\]
Output the interval hull of the accepted set, with \([-1,1]\) as the output if the accepted set is empty:
\[
I^*_{n,\epsilon}=\widehat I^*=
\begin{cases}
\operatorname{hull}\{\vartheta\in[-1,1]:|B_N-\vartheta B_D|\le\rho B_D\},
& \{\vartheta\in[-1,1]:|B_N-\vartheta B_D|\le\rho B_D\}\ne\varnothing,\\
[-1,1],
& \{\vartheta\in[-1,1]:|B_N-\vartheta B_D|\le\rho B_D\}=\varnothing.
\end{cases}
\]
\end{enumerate}
\end{algorithmv}

\begin{theoremv}{thm:uniform-private-upper}[Uniform private upper bounds]*
Let the public parameters satisfy
\begin{itemize}
\item \textbf{(Sample size.)} \(n\in\mathbb N\) and \(n\ge2\).
\item \textbf{(Cell count.)} \(k\in\mathbb N\) and \(k\ge2\).
\item \textbf{(Privacy budget.)} \(0<\epsilon\le1\).
\item \textbf{(Localization radius.)} \(0<h\le1/4\).
\end{itemize}
Define the cell width \(\delta\), public occupancy scale \(\mathcal A\), and public squared-error envelope \(V(h,\delta)\) by
\[
\delta=\frac{2h}{k},
\qquad
\mathcal A=nh\min\{1,n\delta\},
\qquad
V(h,\delta)=2^{20}\left(
h^2+\delta^{2/5}+\mathcal A^{-1}+(\epsilon\mathcal A)^{-2}
\right).
\]
Define the numerical benchmark for public tuning by
\[
r(n,\epsilon)=
\min\left\{1,\,
\max\left\{
n^{-1/4},\,
(n^2\epsilon)^{-1/7},\,
(n\epsilon)^{-1/2}
\right\}\right\}.
\]
The joint release in \cref{obj:def:private-ratio} is
\[
\widetilde S=
\bigl(S_N(D)+\xi_N,\ S_D(D)+\xi_D\bigr),
\qquad
\xi_N,\xi_D\overset{\mathrm{ind}}{\sim}
\operatorname{Laplace}(0,12/\epsilon).
\]
This release and its scalar and interval postprocessings
\(\widehat T_{h,k}\) and \(\widehat I_{h,k}\) in
\cref{obj:def:private-ratio}, together with the publicly tuned outputs
\(\widehat T^*\) and \(\widehat I^*\) in \cref{obj:def:public-tuning},
are everywhere-defined Markov kernels satisfying pure
\(\epsilon\)-replacement privacy on all of \(\mathcal O^n\) in the sense of
\cref{obj:def:private-kernels}.

For every \(P\in\mathcal P\), the causal-law class in
\cref{obj:def:complete-model}, let \(\theta(P)=\tau_P(x_0)\) be the point
target at \(x_0=1/2\). Under the observed sample law
\((P^{\mathrm{obs}})^{\otimes n}\), the local outputs satisfy
\[
E_{(P^{\mathrm{obs}})^{\otimes n},\widehat T_{h,k}}
\left[(\widehat T_{h,k}(D)-\theta(P))^2\right]
\le V(h,\delta),
\]
\[
\Pr_{(P^{\mathrm{obs}})^{\otimes n},\widehat I_{h,k}}
\left\{\theta(P)\in\widehat I_{h,k}(D)\right\}
\ge\frac9{10},
\qquad
E_{(P^{\mathrm{obs}})^{\otimes n},\widehat I_{h,k}}
\operatorname{Leb}(\widehat I_{h,k}(D))
\le\min\left\{2,\ 2\sqrt{10V(h,\delta)}\right\}.
\]
The publicly tuned outputs satisfy
\[
\sup_{P\in\mathcal P}
E_{(P^{\mathrm{obs}})^{\otimes n},\widehat T^*}
\left|\widehat T^*(D)-\theta(P)\right|
\le2^{18}r(n,\epsilon),
\]
\[
\sup_{P\in\mathcal P}
E_{(P^{\mathrm{obs}})^{\otimes n},\widehat I^*}
\operatorname{Leb}(\widehat I^*(D))
\le2^{22}r(n,\epsilon),
\]
and, for every \(P\in\mathcal P\),
\[
\Pr_{(P^{\mathrm{obs}})^{\otimes n},\widehat I^*}
\left\{\theta(P)\in\widehat I^*(D)\right\}
\ge\frac9{10}.
\]
\end{theoremv}
% CAUSALSMITH-CITED-SCOPE-BEGIN thm:uniform-private-upper
\verificationfootnotetext{\textbf{Formalization scope.} This result uses the cited conclusion \citep{BoucheronLugosiMassart2013} (Theorem 3.1, replacement-copy form, printed page 54). The cited source proof is not formalized here: Lean verifies its use and all remaining steps. If that cited conclusion is false, the portions of this result that depend on it are not certified.}
% CAUSALSMITH-CITED-SCOPE-END thm:uniform-private-upper

\begin{theoremv}{thm:matched-risk-frontier}[Matched private risk frontier]*
For every integer \(n\ge2\) and privacy parameter \(0<\epsilon\le1\), define the numerical benchmark
\[
r(n,\epsilon)
=
\min\left\{1,\,
\max\left\{n^{-1/4},\,(n^2\epsilon)^{-1/7},\,(n\epsilon)^{-1/2}\right\}\right\}.
\]
The private minimax absolute risk \(R_{n,\epsilon}\) of \cref{obj:def:risk-criterion}, with the infimum taken over all everywhere-defined randomized pure-private scalar Markov kernels as in \cref{obj:def:private-kernels}, satisfies
\[
2^{-20}r(n,\epsilon)
\le R_{n,\epsilon}
\le 2^{18}r(n,\epsilon).
\]
Moreover, the publicly tuned scalar estimator \(\widehat T^*\), which releases zero when \(r(n,\epsilon)\ge1/8\) and otherwise uses the specified publicly tuned dyadic pair-ratio release, is an everywhere-defined Markov kernel satisfying pure replacement \(\epsilon\)-differential privacy. Its maximal expected absolute error satisfies
\[
\sup_{P\in\mathcal P}
E_{(P^{\mathrm{obs}})^{\otimes n},\widehat T^*}
\left|\widehat T^*(D)-\theta(P)\right|
\le 2^{18}r(n,\epsilon),
\]
where \(\mathcal P\) is the complete causal-law class and \(\theta(P)\) is its pointwise treatment-effect target, with expectation integrating the observed sample and release randomization as in \cref{obj:def:risk-criterion}. The supremum includes every admissible measurable covariate density \(f_P\) in that class.
\end{theoremv}
% CAUSALSMITH-CITED-SCOPE-BEGIN thm:matched-risk-frontier
\verificationfootnotetext{\textbf{Formalization scope.} This result uses the cited conclusion \citep{BoucheronLugosiMassart2013} (Theorem 3.1, replacement-copy form, printed page 54), \citep{BoucheronLugosiMassart2013} (Theorem 6.2, displayed log-mgf conclusion in its proof, printed page 167) and \citep{KellerTrotterAppliedCombinatorics} (Section 5.6, Theorem 5.40). The cited source proof is not formalized here: Lean verifies its use and all remaining steps. If that cited conclusion is false, the portions of this result that depend on it are not certified.}
% CAUSALSMITH-CITED-SCOPE-END thm:matched-risk-frontier

\begin{theoremv}{thm:sharp-interval-frontier}[Sharp honest interval frontier]*
For every integer \(n\ge2\) and privacy budget \(0<\epsilon\le1\), define the numerical benchmark
\[
r(n,\epsilon)
=
\min\left\{1,\,
\max\left\{
n^{-1/4},
(n^2\epsilon)^{-1/7},
(n\epsilon)^{-1/2}
\right\}\right\}.
\]
The honest expected-length benchmark and its constants satisfy
\[
\ell(n,\epsilon)=r(n,\epsilon)>0,
\qquad
c_I=2^{-20}>0,
\qquad
C_I=2^{22}>0.
\]

The interval kernel \(I^*_{n,\epsilon}\) defined by candidate-value inversion in \cref{obj:def:interval-handle} equals the total publicly tuned connected-interval kernel of \cref{obj:def:public-tuning}:
\[
I^*_{n,\epsilon}=\widehat I^*.
\]
For \(r(n,\epsilon)\ge1/8\), this kernel returns \([-1,1]\). For \(r(n,\epsilon)<1/8\), its public tuning is
\[
h=2^{-\lceil\log_2(1/r(n,\epsilon))\rceil},
\qquad
k=2h^{-4},
\qquad
\delta=\frac{2h}{k}=h^5,
\]
where \(0<h\le1/4\) and \(k\) is an integer with \(k\ge2\). Define the public occupancy scale and mean-squared-error envelope by
\[
\mathcal A=nh\min\{1,n\delta\},
\qquad
V(h,\delta)
=
2^{20}\left[
h^2+\delta^{2/5}+\mathcal A^{-1}
+(\epsilon\mathcal A)^{-2}
\right].
\]
Using the clipped pair-ratio estimator from the single joint release in \cref{obj:def:private-ratio}, the tuned interval is
\[
\widehat I^*(D)
=
[-1,1]\cap
\left[
\widehat T_{h,k}(D)-\sqrt{10V(h,\delta)},
\widehat T_{h,k}(D)+\sqrt{10V(h,\delta)}
\right].
\]

The kernel \(I^*_{n,\epsilon}\) is a Markov kernel satisfying pure replacement \(\epsilon\)-differential privacy in the sense of \cref{obj:def:private-kernels}. For the causal-law class \(\mathcal P\) in \cref{obj:def:complete-model},
\[
\Pr_{(P^{\mathrm{obs}})^{\otimes n},I^*_{n,\epsilon}}
\left\{\theta(P)\in I^*_{n,\epsilon}(D)\right\}
\ge\frac9{10}
\qquad\text{for every }P\in\mathcal P.
\]
The honest private expected-length frontier \(H_{n,\epsilon}\) of \cref{obj:def:interval-criterion}, whose infimum ranges over all total private Markov connected-interval kernels with this class-wide coverage, satisfies
\[
2^{-20}r(n,\epsilon)
\le H_{n,\epsilon}
\le
\sup_{P\in\mathcal P}
E_{(P^{\mathrm{obs}})^{\otimes n},I^*_{n,\epsilon}}
\operatorname{Leb}\!\left(I^*_{n,\epsilon}(D)\right)
\le2^{22}r(n,\epsilon).
\]
For the private minimax absolute-error risk \(R_{n,\epsilon}\) of \cref{obj:def:risk-criterion}, the extended nonnegative real frontiers also satisfy
\[
2^{-38}R_{n,\epsilon}
\le H_{n,\epsilon}
\le2^{42}R_{n,\epsilon}.
\]
\end{theoremv}
% CAUSALSMITH-CITED-SCOPE-BEGIN thm:sharp-interval-frontier
\verificationfootnotetext{\textbf{Formalization scope.} This result uses the cited conclusion \citep{BoucheronLugosiMassart2013} (Theorem 3.1, replacement-copy form, printed page 54), \citep{BoucheronLugosiMassart2013} (Theorem 6.2, displayed log-mgf conclusion in its proof, printed page 167) and \citep{KellerTrotterAppliedCombinatorics} (Section 5.6, Theorem 5.40). The cited source proof is not formalized here: Lean verifies its use and all remaining steps. If that cited conclusion is false, the portions of this result that depend on it are not certified.}
% CAUSALSMITH-CITED-SCOPE-END thm:sharp-interval-frontier

\begin{propositionv}{prop:regime-and-sanity}[Benchmark regimes and consistency]*
For every integer \(n\ge2\) and every privacy budget \(0<\epsilon\le1\), define the numerical benchmark
\[
r(n,\epsilon)
=
\min\left\{1,\,
\max\left\{
n^{-1/4},\,
(n^2\epsilon)^{-1/7},\,
(n\epsilon)^{-1/2}
\right\}\right\}.
\]
Its regime identities are
\[
\begin{aligned}
n^{-1/4}\le\epsilon
&\ \Longrightarrow\ r(n,\epsilon)=n^{-1/4},\\
n^{-3/5}\le\epsilon\le n^{-1/4}
&\ \Longrightarrow\ r(n,\epsilon)=(n^2\epsilon)^{-1/7},\\
n^{-1}\le\epsilon\le n^{-3/5}
&\ \Longrightarrow\ r(n,\epsilon)=(n\epsilon)^{-1/2},\\
\epsilon\le n^{-1}
&\ \Longrightarrow\ r(n,\epsilon)=1.
\end{aligned}
\]
In particular,
\[
r(n,1)=n^{-1/4}.
\]

For every budget sequence \((\epsilon_n)_{n\in\mathbb N}\) satisfying \(0<\epsilon_n\le1\) for every \(n\ge2\), the private minimax absolute risk \(R_{n,\epsilon_n}\) defined in \cref{obj:def:risk-criterion}, with values in \([0,\infty]\), satisfies
\[
R_{n,\epsilon_n}\longrightarrow0
\quad\Longleftrightarrow\quad
n\epsilon_n\longrightarrow+\infty
\qquad(n\to\infty).
\]
\end{propositionv}
% CAUSALSMITH-CITED-SCOPE-BEGIN prop:regime-and-sanity
\verificationfootnotetext{\textbf{Formalization scope.} This result uses the cited conclusion \citep{BoucheronLugosiMassart2013} (Theorem 3.1, replacement-copy form, printed page 54), \citep{BoucheronLugosiMassart2013} (Theorem 6.2, displayed log-mgf conclusion in its proof, printed page 167) and \citep{KellerTrotterAppliedCombinatorics} (Section 5.6, Theorem 5.40). The cited source proof is not formalized here: Lean verifies its use and all remaining steps. If that cited conclusion is false, the portions of this result that depend on it are not certified.}
% CAUSALSMITH-CITED-SCOPE-END prop:regime-and-sanity

\begin{lemmav}{lem:point-version}[Continuous contrast and point identification]
Let \(P\in\mathcal P\), the class of causal laws specified in \cref{obj:def:complete-model}. Write \(P_X\) for the law of the covariate \(X\), and \(P^{\mathrm{obs}}\) for the law of the observed record \((X,A,Y)\). The selected observational contrast \(\tau_P\) is the difference of the selected treated and control outcome regressions:
\[
\tau_P(x)
=
E_P[Y\mid X=x,A=1]
-
E_P[Y\mid X=x,A=0],
\qquad x\in[0,1].
\]
Here the conditional means denote the selected regression versions carried by \(P\). Define the point target by
\[
x_0=\tfrac12,
\qquad
\theta(P)=\tau_P(x_0).
\]

Then \(\tau_P\) is continuous on \([0,1]\). For every continuous function \(v:[0,1]\to\mathbb R\),
\[
v=\tau_P\quad P_X\text{-almost everywhere}
\quad\Longrightarrow\quad
v(x)=\tau_P(x)\quad\text{for every }x\in[0,1].
\]
Moreover, the point target and the potential-outcome contrast satisfy
\[
\theta(P)\in[-1,1],
\qquad
\tau_P(X)=E_P[Y(1)-Y(0)\mid X]
\quad P\text{-almost surely},
\]
where \(Y(0)\) and \(Y(1)\) are the binary control and treated potential outcomes. Finally, for every \(P'\in\mathcal P\),
\[
P'^{\mathrm{obs}}=P^{\mathrm{obs}}
\quad\Longrightarrow\quad
\theta(P')=\theta(P).
\]
\end{lemmav}

\begin{lemmav}{lem:occupancy-cancellation}[Occupancy cancellation and ratio bias]
Suppose the following conditions hold:
\begin{itemize}
\item \textbf{(Sample and cells.)} The sample size \(n\) and public cell count \(k\) are integers satisfying \(n\ge 2\) and \(k\ge 2\).
\item \textbf{(Localization.)} The public localization radius \(h\) satisfies \(0<h\le 1/4\).
\item \textbf{(Model.)} The causal law \(P\) belongs to the class \(\mathcal P\) of \cref{obj:def:complete-model}.
\item \textbf{(Sampling.)} The dataset law \(Q\) satisfies \cref{obj:ass:iid} for sample size \(n\) and causal law \(P\).
\end{itemize}
Write \(\delta=2h/k\) for the cell width, and define the public bias envelope \(b\) and occupancy scale \(\mathcal A\) by
\[
b=3h+24\delta^{1/5},
\qquad
\mathcal A=nh\min\{1,n\delta\}.
\]
Let \(\overline N\) and \(\overline D\) be the population means of the numerator and denominator statistics:
\[
\overline N=\int S_N(D)\,dQ(D),
\qquad
\overline D=\int S_D(D)\,dQ(D).
\]
For each public cell \(j=1,\ldots,k\), let \(\nu_j\) and \(d_j\) be the expectations of the numerator and denominator pair contributions, respectively, under two independent observations drawn from the conditional observed law in that cell. Define its occupancy weight by
\[
w_j
=
nP_X(\text{cell }j)
\left[1-\left(1-P_X(\text{cell }j)\right)^{n-1}\right].
\]
With \(W(P)\) denoting the total population occupancy weight, \(\theta(P)=\tau_P(1/2)\) the point target, and \(d_0\) the public denominator floor, the population moments satisfy
\[
\overline N=\sum_{j=1}^{k}w_j\nu_j,
\qquad
\overline D=\sum_{j=1}^{k}w_jd_j.
\]
For every \(j=1,\ldots,k\),
\[
d_j\ge\frac38,
\qquad
\left|\nu_j-\theta(P)d_j\right|\le b\,d_j.
\]
Moreover,
\[
\frac{\mathcal A}{8}\le W(P)\le\frac{9\mathcal A}{2},
\qquad
\overline D\ge d_0,
\qquad
\left|\frac{\overline N}{\overline D}-\theta(P)\right|\le b.
\]
\end{lemmav}

\begin{lemmav}{lem:all-count-stability}[All-count stability and variance]*
Let \(S_N(D)\) and \(S_D(D)\) be the sums over the \(k\) public cells of the count-weighted within-cell treatment–outcome pair contributions and treatment pair contributions, respectively. Let \(W(P)\), the total population occupancy weight, be
\[
W(P)=\sum_{j=1}^{k}w_j,
\]
where \(w_j\) is the population occupancy weight of cell \(j\). Suppose:
\begin{itemize}
\item \textbf{(Parameters.)} The integers \(n,k\) satisfy \(n\ge2\) and \(k\ge2\), and the public window radius \(h\) satisfies \(0<h\le1/4\).
\item \textbf{(Model.)} The causal law \(P\) belongs to \(\mathcal P\) as specified in \cref{obj:def:complete-model}.
\item \textbf{(Sampling.)} The dataset law satisfies \cref{obj:ass:iid}.
\end{itemize}
For every pair of ordered datasets \(D,D^{\prime}\in\mathcal O^n\) whose replacement Hamming distance, the number of differing record positions, is \(d_{\mathrm H}(D,D^{\prime})=1\),
\[
\left|S_N(D)-S_N(D^{\prime})\right|
+\left|S_D(D)-S_D(D^{\prime})\right|
\le12.
\]
Under the stated dataset law,
\[
\max\left\{
\operatorname{Var}(S_N(D)),
\operatorname{Var}(S_D(D))
\right\}
\le18W(P).
\]
The replacement bound applies to every input dataset, including cells with occupancies zero, one, and two.
\end{lemmav}
% CAUSALSMITH-CITED-SCOPE-BEGIN lem:all-count-stability
\verificationfootnotetext{\textbf{Formalization scope.} This result uses the cited conclusion \citep{BoucheronLugosiMassart2013} (Theorem 3.1, replacement-copy form, printed page 54). The cited source proof is not formalized here: Lean verifies its use and all remaining steps. If that cited conclusion is false, the portions of this result that depend on it are not certified.}
% CAUSALSMITH-CITED-SCOPE-END lem:all-count-stability

\begin{definitionv}{synth_6}[Binary sign vector encoding]
The sign vector \(\lambda\) is a binary encoding of signs in \(\{-1,1\}\), with one coordinate for each active integer index determined by the location \(x_0\), macro radius \(h_{\mathrm L}\), and micro radius \(\delta_{\mathrm L}\). Specifically, \(\lambda\) is any function
\[
\lambda:
\left\{
j\in\mathbb Z:
\begin{gathered}
\left\lfloor\frac{x_0-h_{\mathrm L}}{\delta_{\mathrm L}}\right\rfloor-1
\le j\le
\left\lceil\frac{x_0+h_{\mathrm L}}{\delta_{\mathrm L}}\right\rceil+1,\\
j\delta_{\mathrm L}-\delta_{\mathrm L}<x_0+h_{\mathrm L},
\qquad
x_0-h_{\mathrm L}<j\delta_{\mathrm L}+\delta_{\mathrm L}
\end{gathered}
\right\}
\longrightarrow \{0,1\}.
\]
\end{definitionv}

\begin{definitionv}{synth_5}[Target separation amplitude]
The target-separation amplitude \(t\) is defined by
\[
t := 2^{-12}h_{\mathrm L},
\]
where \(h_{\mathrm L}\) is the macro localization radius for the lower-bound experiments.
\end{definitionv}

\begin{definitionv}{def:cosine-family}[Shared-sign causal alternative family]
The causal alternatives \(P_\lambda\) and their observed-sample mixture \(Q_n\) are defined with localization center \(x_0=1/2\), macro localization radius \(0<h_{\mathrm L}\le 1/4\), micro localization radius \(\delta_{\mathrm L}\), and target-separation amplitude \(t\), where
\[
\delta_{\mathrm L}=h_{\mathrm L}^{5},
\qquad
t=2^{-12}h_{\mathrm L}.
\]
Each sign vector \(\lambda\) assigns a binary value \(\lambda_j\in\{0,1\}\) to every active integer index in
\[
\left\{
j\in\mathbb Z:
\left\lfloor\frac{x_0-h_{\mathrm L}}{\delta_{\mathrm L}}\right\rfloor-1
\le j\le
\left\lceil\frac{x_0+h_{\mathrm L}}{\delta_{\mathrm L}}\right\rceil+1,\quad
j\delta_{\mathrm L}-\delta_{\mathrm L}<x_0+h_{\mathrm L},\quad
x_0-h_{\mathrm L}<j\delta_{\mathrm L}+\delta_{\mathrm L}
\right\}.
\]
The corresponding signed coefficient is \(2\lambda_j-1\in\{-1,1\}\). The localized contrast profile \(q(x)\) and shared-sign perturbation \(S_\lambda(x)\) are
\[
q(x)=
\begin{cases}
\displaystyle
\cos^{4}\!\left(\frac{\pi(x-x_0)}{2h_{\mathrm L}}\right),
& |x-x_0|\le h_{\mathrm L},\\
0,& |x-x_0|>h_{\mathrm L},
\end{cases}
\]
and
\[
S_\lambda(x)=
\begin{cases}
\displaystyle
\cos^{2}\!\left(\frac{\pi(x-x_0)}{2h_{\mathrm L}}\right)
\sum_j(2\lambda_j-1)
\begin{cases}
\displaystyle
\cos\!\left(\frac{\pi(x-j\delta_{\mathrm L})}{2\delta_{\mathrm L}}\right),
& |x-j\delta_{\mathrm L}|\le\delta_{\mathrm L},\\
0,& |x-j\delta_{\mathrm L}|>\delta_{\mathrm L},
\end{cases}
& |x-x_0|\le h_{\mathrm L},\\
0,& |x-x_0|>h_{\mathrm L},
\end{cases}
\]
where the sum runs over the active indices above.

For each \(\lambda\), the causal law \(P_\lambda\) has covariate \(X\sim\operatorname{Unif}[0,1]\). Conditional on \(X=x\), the binary treatment indicator \(A\) and binary potential outcomes \(Y(0),Y(1)\) are mutually independent, with
\[
\begin{aligned}
A\mid X=x
&\sim\operatorname{Bernoulli}
\left(\frac{1+\sqrt t\,S_\lambda(x)}{2}\right),\\
Y(0)\mid X=x
&\sim\operatorname{Bernoulli}
\left(\frac{1-\sqrt t\,S_\lambda(x)-tq(x)}{2}\right),\\
Y(1)\mid X=x
&\sim\operatorname{Bernoulli}
\left(\frac{1-\sqrt t\,S_\lambda(x)+tq(x)}{2}\right).
\end{aligned}
\]
The observed outcome is \(Y=Y(A)\). The selected covariate density is \(f_{P_\lambda}(x)=1\), and the displayed Bernoulli parameters specify the selected propensity and arm-mean versions.

For every nonnegative integer \(n\), let \(P_\lambda^{\mathrm{obs}}\) be the marginal law of \((X,A,Y)\). Define
\[
Q_n=E_\lambda\!\left[(P_\lambda^{\mathrm{obs}})^{\otimes n}\right],
\]
where \(E_\lambda\) is uniform averaging over all binary sign vectors on the active index set, equivalently assigning independent fair binary values to its indices.
\end{definitionv}

\begin{lemmav}{lem:positive-family-certificate}[Positive family and mixture certificate]*
For every \(h_{\mathrm L}\in(0,1/4]\), the macro localization radius, and every integer \(n\ge0\), set the target-separation amplitude \(t\) and micro localization radius \(\delta_{\mathrm L}\) to
\[
t=2^{-12}h_{\mathrm L},
\qquad
\delta_{\mathrm L}=h_{\mathrm L}^{5}.
\]
Let \(P_0\) be the causal null with a uniform covariate and conditionally independent fair binary treatment and potential outcomes, with the observed outcome determined by consistency. Use the alternatives \(P_\lambda\), localized profile \(q(x)\), nuisance perturbation \(S_\lambda(x)\), uniform sign average \(E_\lambda\), and mixture \(Q_n\) from \cref{obj:def:cosine-family}. Throughout, \(\lambda\) is the binary encoding on the active frame indices:
\[
\lambda_j\in\{0,1\},
\qquad
2\lambda_j-1\in\{-1,1\}.
\]
Thus the coefficients of the cosine-frame sum defining \(S_\lambda(x)\) are \(2\lambda_j-1\).

For the model class \(\mathcal P\) in \cref{obj:def:complete-model} and the point target \(\theta(P)=\tau_P(1/2)\),
\[
P_0\in\mathcal P,\qquad
\theta(P_0)=0,\qquad
P_\lambda\in\mathcal P,\qquad
\theta(P_\lambda)=t
\quad\text{for every }\lambda.
\]
For every sign vector \(\lambda\), covariate value \(x\), and binary treatment and outcome values \(A,Y\), define the likelihood marks
\[
U=2A-1,\qquad V_Y=2Y-1.
\]
The conditional likelihood \(L_\lambda(x,U,V_Y)\), equal to four times the alternative conditional probability of the cell \((A,Y)\) given \(x\), satisfies
\[
\begin{aligned}
L_\lambda(x,U,V_Y)
={}&1+U\sqrt t\,S_\lambda(x)
 +V_Y\sqrt t\,\bigl(tq(x)-1\bigr)S_\lambda(x)\\
&+UV_Yt\bigl(q(x)-S_\lambda(x)^2\bigr).
\end{aligned}
\]
The one-record mixture equals the observed null law:
\[
E_\lambda P_\lambda^{\mathrm{obs}}=P_0^{\mathrm{obs}}.
\]
Moreover,
\[
\chi^2\!\left(Q_n,(P_0^{\mathrm{obs}})^{\otimes n}\right)
\le
\exp\!\left(160n^2t^2h_{\mathrm L}\delta_{\mathrm L}\right)-1,
\qquad
Q_n\ll(P_0^{\mathrm{obs}})^{\otimes n},
\]
and the centered likelihood ratio is square-integrable:
\[
\int_{\mathcal O^n}
\left(
\frac{dQ_n}{d(P_0^{\mathrm{obs}})^{\otimes n}}-1
\right)^2
\,d(P_0^{\mathrm{obs}})^{\otimes n}
<\infty.
\]
\end{lemmav}
% CAUSALSMITH-CITED-SCOPE-BEGIN lem:positive-family-certificate
\verificationfootnotetext{\textbf{Formalization scope.} This result uses the cited conclusion \citep{BoucheronLugosiMassart2013} (Theorem 6.2, displayed log-mgf conclusion in its proof, printed page 167). The cited source proof is not formalized here: Lean verifies its use and all remaining steps. If that cited conclusion is false, the portions of this result that depend on it are not certified.}
% CAUSALSMITH-CITED-SCOPE-END lem:positive-family-certificate

\begin{definitionv}{synth_14}[Conditional shared-sign graph]
Let \(x_0=1/2\), \(0<h_{\mathrm L}\le1/4\), and \(\delta_{\mathrm L}=h_{\mathrm L}^5\). Define the macro envelope and cosine frame by
\[
g(x)=\cos^2\!\left\{\frac{\pi(x-x_0)}{2h_{\mathrm L}}\right\}
\mathbf1\{|x-x_0|\le h_{\mathrm L}\},
\qquad
\phi_j(x)=\cos\!\left\{\frac{\pi(x-j\delta_{\mathrm L})}{2\delta_{\mathrm L}}\right\}
\mathbf1\{|x-j\delta_{\mathrm L}|\le\delta_{\mathrm L}\},
\]
and define the finite active-sign index set
\[
\mathcal J=
\left\{j\in\mathbb Z:
(j\delta_{\mathrm L}-\delta_{\mathrm L},j\delta_{\mathrm L}+\delta_{\mathrm L})
\cap(x_0-h_{\mathrm L},x_0+h_{\mathrm L})\ne\varnothing
\right\}.
\]
For a covariate vector \(X=(X_1,\ldots,X_n)\), with \(n\ge0\), the conditional shared-sign graph \(G_X\) is the simple undirected graph with vertex set \(\{1,\ldots,n\}\). Distinct vertices \(i,l\) are adjacent precisely when
\[
\exists j\in\mathcal J:\quad
g(X_i)\phi_j(X_i)\ne0
\quad\text{and}\quad
g(X_l)\phi_j(X_l)\ne0.
\]
Its connected components are selected in lexicographic vertex order.
\end{definitionv}

\begin{definitionv}{synth_7}[Total common component mass]
The total common mass \(\alpha\), for fixed covariate values in \([0,1]\) and a finite vertex set \(C\), is defined by
\[
\alpha
=
\sum_{z:C\to\{0,1\}^2}
\min\bigl\{Q_C(\{z\}),\,P_{0,C}(\{z\})\bigr\},
\]
where \(z\) assigns a binary treatment–outcome pair to each vertex in \(C\), and \(Q_C\) and \(P_{0,C}\) are the alternative and null component mark laws conditional on the fixed covariate values.
\end{definitionv}

\begin{definitionv}{synth_15}[Component mark masses and common overlap]
Fix the covariates and a connected component \(C\) of \(G_X\), and write \(m=|C|\). Order its vertices increasingly and encode their treatment–outcome marks \((A_i,Y_i)_{i\in C}\) as elements of \(\{0,1\}^{2m}\). Let \(Q_C\) be their conditional law under the alternative dataset mixture
\[
Q_n=E_\lambda[(P_\lambda^{\mathrm{obs}})^{\otimes n}]
\]
from \(\cref{obj:def:cosine-family}\). Let \(P_{0,C}\) be their conditional law under \((P_0^{\mathrm{obs}})^{\otimes n}\), where the null observed law has \(X\sim\operatorname{Unif}[0,1]\) and, conditional on \(X\), independent fair Bernoulli treatment and outcome marks. Thus \(P_{0,C}\) is uniform on \(\{0,1\}^{2m}\). For every configuration \(v\in\{0,1\}^{2m}\), define
\[
a_v=Q_C(\{v\}),\qquad
b_v=P_{0,C}(\{v\}),\qquad
c_v=\min\{Q_C(\{v\}),P_{0,C}(\{v\})\}.
\]
In particular, \(a_z=Q_C(\{z\})\), \(b_w=P_{0,C}(\{w\})\), and \(c_z\) is the common diagonal mass at \(z\). Define the total common mass by
\[
\alpha=\sum_{v\in\{0,1\}^{2m}}c_v.
\]
\end{definitionv}

\begin{definitionv}{def:common-mass-coupling}[Component coupling \(\Gamma_C\)]
For component mark configurations \(z,w\in\{0,1\}^{2m}\), define
\[
\Gamma_C(z,w)
=
\begin{cases}
\displaystyle
c_z\mathbf 1\{z=w\}
+\frac{(a_z-c_z)(b_w-c_w)}{1-\alpha},
&\alpha<1,\\[6pt]
c_z\mathbf 1\{z=w\},
&\alpha=1.
\end{cases}
\]
The dataset coupling uses the same covariates in both coordinates and the product of these kernels across the lexicographically selected graph components.
\end{definitionv}

\begin{lemmav}{lem:measurable-component-contraction}[Measurable component coupling and contraction]*
Let \(n\) be a nonnegative integer. Impose the following conditions:
\begin{itemize}
\item \textbf{(Lower-family calibration.)} The macro localization radius \(h_{\mathrm L}\), micro localization radius \(\delta_{\mathrm L}\), and target-separation amplitude \(t\) satisfy
\[
0<h_{\mathrm L}\le\frac14,
\qquad
\delta_{\mathrm L}=h_{\mathrm L}^{5},
\qquad
t=2^{-12}h_{\mathrm L}.
\]
\item \textbf{(Output space.)} Let \(\mathcal B\) be a standard Borel output space.
\item \textbf{(Private release.)} Let \(0<\epsilon\le1\) and let \(M\in\mathfrak M_{n,\epsilon}(\mathcal B)\), as defined in \cref{obj:def:private-kernels}.
\end{itemize}
Let \(Q_n\) be the uniform sign mixture of the alternative observed product laws in \cref{obj:def:cosine-family}, and let \((P_0^{\mathrm{obs}})^{\otimes n}\) be the fair-null dataset law.

For every covariate vector \(X=x\), let \(G_X\) be the shared-sign graph, whose connected components group records linked by shared latent signs. For each connected component \(C\), set \(m=|C|\). Its conditional alternative and null mark laws \(Q_C\) and \(P_{0,C}\) have masses
\[
Q_C(\{z\})
 = E_\lambda\prod_{i\in C}
   \frac{L_\lambda(x_i,U_i,V_{Y,i})}{4},
\qquad
P_{0,C}(\{z\})=4^{-m},
\]
where \(z=((U_i,V_{Y,i}))_{i\in C}\) is a configuration of the binary treatment and outcome marks, \(L_\lambda\) is their conditional likelihood relative to the fair null, and \(E_\lambda\) averages uniformly over the latent sign vectors.

For every \(x\) and every full mark configuration \(z=((U_i,V_{Y,i}))_{i=1}^n\), the conditional masses factor over the components of \(G_X\):
\[
E_\lambda\prod_{i=1}^n
  \frac{L_\lambda(x_i,U_i,V_{Y,i})}{4}
 = \prod_{C\text{ component of }G_X} Q_C(\{z|_C\}),
\qquad
4^{-n}
 = \prod_{C\text{ component of }G_X} P_{0,C}(\{z|_C\}).
\]
For every such component,
\[
m=1\ \Longrightarrow\ Q_C=P_{0,C},
\qquad
d_{\mathrm{TV}}(Q_C,P_{0,C})\le4tm^2,
\]
where \(d_{\mathrm{TV}}\) denotes total variation distance. The common-mass measure \(\Gamma_C\) of \cref{obj:def:common-mass-coupling} is a coupling of \(Q_C\) and \(P_{0,C}\). The product of these component couplings, with components ordered by increasing least record index, defines a measurable conditional dataset coupling. Integrating it over the shared covariate law, the product of \(n\) uniform laws on \([0,1]\), gives a dataset coupling with marginals \(Q_n\) and \((P_0^{\mathrm{obs}})^{\otimes n}\).

If \(n\delta_{\mathrm L}\le1/128\), then
\[
W_{\mathrm H}\!\left(Q_n,(P_0^{\mathrm{obs}})^{\otimes n}\right)
 \le1024tn^2h_{\mathrm L}\delta_{\mathrm L},
\qquad
d_{\mathrm{TV}}\!\left(MQ_n,M(P_0^{\mathrm{obs}})^{\otimes n}\right)
 \le2048\epsilon tn^2h_{\mathrm L}\delta_{\mathrm L},
\]
where \(W_{\mathrm H}\) is the infimum of expected dataset Hamming distance over couplings, Hamming distance counts the record positions at which two datasets differ, and \(MQ\) denotes the output law obtained by applying \(M\) to a dataset with law \(Q\).
\end{lemmav}
% CAUSALSMITH-CITED-SCOPE-BEGIN lem:measurable-component-contraction
\verificationfootnotetext{\textbf{Formalization scope.} This result uses the cited conclusion \citep{KellerTrotterAppliedCombinatorics} (Section 5.6, Theorem 5.40) and \citep{BoucheronLugosiMassart2013} (Theorem 6.2, displayed log-mgf conclusion in its proof, printed page 167). The cited source proof is not formalized here: Lean verifies its use and all remaining steps. If that cited conclusion is false, the portions of this result that depend on it are not certified.}
% CAUSALSMITH-CITED-SCOPE-END lem:measurable-component-contraction

\begin{definitionv}{def:direct-family}[Direct localized causal alternative]
The causal law \(P_1\), the direct localized alternative, is defined for the macro localization radius \(0<h_{\mathrm L}\le 1/4\), with target-separation amplitude and localized contrast profile
\[
t=2^{-12}h_{\mathrm L},
\qquad
q(x)=
\begin{cases}
\cos^4\!\left(\dfrac{\pi(x-x_0)}{2h_{\mathrm L}}\right),
& |x-x_0|\le h_{\mathrm L},\\
0,& |x-x_0|>h_{\mathrm L},
\end{cases}
\quad x\in[0,1],
\]
where \(x_0\) is the target covariate location. Under \(P_1\), draw the covariate \(X\sim\operatorname{Unif}[0,1]\). Conditional on \(X\), draw the treatment indicator \(A\) and the potential outcomes \(Y(0)\) and \(Y(1)\) independently according to
\[
A\sim\operatorname{Bernoulli}(1/2),
\qquad
Y(0)\sim\operatorname{Bernoulli}(1/2),
\qquad
Y(1)\sim\operatorname{Bernoulli}\bigl(1/2+tq(X)\bigr).
\]
Set the observed outcome \(Y\) to
\[
Y=Y(A).
\]
The selected covariate density, propensity, control regression, and treatment-effect contrast are
\[
f_{P_1}(x)=1,\qquad
e_{P_1}(x)=\frac12,\qquad
\mu_{0,P_1}(x)=\frac12,\qquad
\tau_{P_1}(x)=tq(x).
\]
\end{definitionv}

\begin{lemmav}{lem:private-kernel-comparison}[Private kernel comparison bounds]
Let \(n\) be a nonnegative integer and let \(Q,Q'\) be probability laws on \(\mathcal O^n\), the space of ordered datasets of \(n\) observed records. Suppose that:
\begin{itemize}
\item \textbf{(Privacy parameter.)} \(0<\epsilon\le 1\).
\item \textbf{(Output space.)} \(\mathcal B\) is a standard Borel output space.
\item \textbf{(Private release.)} \(M\in\mathfrak M_{n,\epsilon}(\mathcal B)\), the class of everywhere probability-valued private Markov kernels defined in \cref{obj:def:private-kernels}.
\end{itemize}
Write \(MQ\) for the induced output law, so that
\[
(MQ)(E)=\int M(D,E)\,Q(dD),
\qquad E\in\mathcal B(\mathcal B).
\]
For probability laws on a common measurable space, let \(d_{\mathrm{TV}}\) denote total variation distance:
\[
d_{\mathrm{TV}}(Q,Q')=\sup_{E\ \mathrm{measurable}}|Q(E)-Q'(E)|.
\]
Let \(W_{\mathrm H}\) denote the infimum of expected Hamming distance over all joint laws with the specified marginals:
\[
W_{\mathrm H}(Q,Q')
=
\inf_{\substack{\mathcal L(D)=Q\\\mathcal L(D')=Q'}}
E\!\left[d_{\mathrm H}(D,D')\right],
\qquad
d_{\mathrm H}(D,D')
=
\sum_{i=1}^{n}\mathbf 1\{D_i\ne D'_i\}.
\]
Then
\[
d_{\mathrm{TV}}(MQ,MQ')
\le 2\epsilon W_{\mathrm H}(Q,Q').
\]
Furthermore, for every everywhere probability-valued Markov kernel \(M\) from \(\mathcal O^n\) to \(\mathcal B\),
\[
d_{\mathrm{TV}}(MQ,MQ')\le d_{\mathrm{TV}}(Q,Q').
\]
If \(Q\ll Q'\) and the squared Radon--Nikodym deviation is integrable, that is,
\[
\int_{\mathcal O^n}
\left(\frac{dQ}{dQ'}(D)-1\right)^2\,Q'(dD)<\infty,
\]
then, with chi-square divergence defined by
\[
\chi^2(Q,Q')
=
\int_{\mathcal O^n}
\left(\frac{dQ}{dQ'}(D)-1\right)^2\,Q'(dD),
\]
one has
\[
d_{\mathrm{TV}}(Q,Q')
\le \frac12\sqrt{\chi^2(Q,Q')}.
\]
\end{lemmav}

\begin{lemmav}{lem:frontier-testing-priors}[Finite testing priors]*
For every integer \(n\ge2\) and privacy budget \(0<\epsilon\le1\), define the numerical benchmark
\[
r(n,\epsilon)
=
\min\left\{1,\max\left\{n^{-1/4},
\max\left\{(n^2\epsilon)^{-1/7},(n\epsilon)^{-1/2}\right\}\right\}\right\}.
\]
There exist two nonempty finite families of laws in \(\mathcal P\), the model class of \cref{obj:def:complete-model}, whose uniformly averaged observed \(n\)-record product laws are \(Q^{(0)}\) and \(Q^{(1)}\), respectively, and a common target separation \(\Delta\) satisfying
\[
\Delta\ge2^{-14}r(n,\epsilon).
\]
Every member of the first family has point target \(\theta(P)=0\), and every member of the second has \(\theta(P)=\Delta\).

These families can be chosen with macro localization radius
\[
h_{\mathrm L}=\frac{r(n,\epsilon)}4,
\qquad 0<h_{\mathrm L}\le\frac14,
\]
and the following properties:
\begin{itemize}
\item \textbf{(Null family.)} Every member of the first family equals \(P_0\), the causal law with uniform covariate \(X\), conditionally independent fair binary treatment and potential outcomes given \(X\), and observed outcome \(Y=Y(A)\). Consequently,
\[
Q^{(0)}=(P_0^{\mathrm{obs}})^{\otimes n}.
\]
\item \textbf{(Sign family.)} If
\[
1<n\epsilon
\quad\text{and}\quad
\left[
\left((n^2\epsilon)^{-1/7}\le n^{-1/4}
\ \text{and}\ (n\epsilon)^{-1/2}\le n^{-1/4}\right)
\ \text{or}\
\left(n^{-1/4}<(n^2\epsilon)^{-1/7}
\ \text{and}\ (n\epsilon)^{-1/2}<(n^2\epsilon)^{-1/7}\right)
\right],
\]
every member of the second family is \(P_\lambda\) for some sign vector \(\lambda\), using the construction of \cref{obj:def:cosine-family} at radius \(h_{\mathrm L}\), and
\[
Q^{(1)}=Q_n,
\]
where \(Q_n\) is the uniform sign mixture at that radius defined in \cref{obj:def:cosine-family}.
\item \textbf{(Direct family.)} If the preceding condition fails, every member of the second family equals \(P_1\), the direct alternative of \cref{obj:def:direct-family} at radius \(h_{\mathrm L}\), and
\[
Q^{(1)}=(P_1^{\mathrm{obs}})^{\otimes n}.
\]
\end{itemize}
For every standard Borel output space \(\mathcal B\) and every private Markov kernel \(M\in\mathfrak M_{n,\epsilon}(\mathcal B)\) as defined in \cref{obj:def:private-kernels}, the induced output laws satisfy
\[
d_{\mathrm{TV}}\!\left(MQ^{(0)},MQ^{(1)}\right)\le\frac14,
\]
where \(d_{\mathrm{TV}}\) denotes total variation distance and \(MQ^{(0)},MQ^{(1)}\) are obtained by applying \(M\) to the respective dataset laws.
\end{lemmav}
% CAUSALSMITH-CITED-SCOPE-BEGIN lem:frontier-testing-priors
\verificationfootnotetext{\textbf{Formalization scope.} This result uses the cited conclusion \citep{BoucheronLugosiMassart2013} (Theorem 6.2, displayed log-mgf conclusion in its proof, printed page 167) and \citep{KellerTrotterAppliedCombinatorics} (Section 5.6, Theorem 5.40). The cited source proof is not formalized here: Lean verifies its use and all remaining steps. If that cited conclusion is false, the portions of this result that depend on it are not certified.}
% CAUSALSMITH-CITED-SCOPE-END lem:frontier-testing-priors
